\documentclass[11pt]{amsart}
\usepackage[T1]{fontenc}
\usepackage{lmodern}
\usepackage{amsmath,amssymb,amsthm}
\usepackage{xcolor}
\usepackage[colorlinks=true,linkcolor=blue!55!black,citecolor=blue!55!black,urlcolor=blue!55!black]{hyperref}
\newtheorem{theorem}{Theorem}[section]
\newtheorem{corollary}[theorem]{Corollary}
\theoremstyle{remark}\newtheorem{remark}[theorem]{Remark}
\title{Nonexistence of a Rational Antipodal 34-Point 7-Design for the Chebyshev Measure}
\author{Zijian Zeng}
\address{Institute of Computer Science and Digital Innovation, UCSI University, Kuala Lumpur, 56000, MALAYSIA}
\email{zijianzeng@foxmail.com; 1002266693@ucsiuniversity.edu.my}
\date{August 31, 2026}
\hypersetup{pdftitle={Nonexistence of a Rational Antipodal 34-Point 7-Design for the Chebyshev Measure},pdfauthor={Zijian Zeng}}
\begin{document}
\begin{abstract}
Mishima, Lu, Sawa, and Uchida proved that rational antipodal $7$-designs
for the Chebyshev measure do not exist with $2N$ points for
$1\leq N\leq16$, and posed the existence problem starting at $N=17$.
We settle the first case negatively. In fact, the three required rational
moment equations in seventeen variables have no solution, even without the
distinctness and interval conditions. The proof is an elementary $2$-adic
argument using a primitive common denominator.
\end{abstract}
\maketitle
\section{Introduction and formulation}
Let $d\mu(t)=dt/(\pi\sqrt{1-t^2})$ be the Chebyshev probability measure on
$(-1,1)$. An antipodal rational $7$-design with $34$ points would be a set
$X=\{\pm x_i:1\leq i\leq17\}$ of distinct rational points in $(-1,1)$
such that
\[
 \frac1{34}\sum_{x\in X}f(x)=\int_{-1}^{1}f(t)\,d\mu(t)
 \qquad (\deg f\leq7).
\]
The even moments are $\int t^{2k}\,d\mu(t)=\binom{2k}{k}/4^k$.
Antipodality cancels the odd moments, so the condition is equivalent to
\begin{equation}\label{eq:moments}
 \sum_{i=1}^{17}x_i^2=\frac{17}{2},\qquad
 \sum_{i=1}^{17}x_i^4=\frac{51}{8},\qquad
 \sum_{i=1}^{17}x_i^6=\frac{85}{16}.
\end{equation}
This equivalence and the prior nonexistence result appear in
\cite[Section~8, preceding Theorem~8.4; Theorem~8.4 and Problem~8.5]{Mishima2026}.
\section{The obstruction}
\begin{theorem}\label{thm:main}
There is no tuple $(x_1,\ldots,x_{17})\in\mathbb Q^{17}$ satisfying
\eqref{eq:moments}.
\end{theorem}
\begin{proof}
Choose a positive common denominator $q$, minimal so that $x_i=a_i/q$ with
$a_i\in\mathbb Z$. Thus $\gcd(q,a_1,\ldots,a_{17})=1$. Clearing
denominators gives
\begin{equation}\label{eq:cleared}
 2\sum a_i^2=17q^2,\qquad 8\sum a_i^4=51q^4,\qquad
 16\sum a_i^6=85q^6.
\end{equation}
The first equality makes $q$ even. Put $s=v_2(q)\geq1$. Minimality ensures
that at least one $a_i$ is odd.

First suppose $s=1$ and write $q=2Q$, where $Q$ is odd. The last two
equalities yield
\[
 \sum a_i^4=102Q^4\equiv2\pmod4,
 \qquad \sum a_i^6=340Q^6\equiv4\pmod8.
\]
If $r$ is the number of odd $a_i$, reduction gives
$r\equiv2\pmod4$ and $r\equiv4\pmod8$, which is impossible.

Now suppose $s\geq2$. The fourth-power equality gives
$v_2(\sum a_i^4)=4s-3\geq5$. Every odd fourth power is $1$ modulo $16$,
and every even fourth power is $0$ modulo $16$. Hence
$r\equiv0\pmod{16}$. Since $1\leq r\leq17$, we have $r=16$; denote the
unique even numerator by $b$.

The square equality gives $8\mid\sum a_i^2$. The sixteen odd squares
contribute $16\equiv0\pmod8$, so $b^2\equiv0\pmod8$ and $4\mid b$.
For the odd numerators write $a_i^2=1+8u_i$, and set $U=\sum u_i$.
The terms $b^4$ and $b^6$ vanish modulo $64$. Since the sixth-power sum is
divisible by $2^{6s-4}$, it vanishes modulo $64$, and
\[
 0\equiv\sum a_i^6\equiv16+24U\pmod{64}.
\]
Thus $U\equiv2\pmod8$. But then
\[
 \sum a_i^4\equiv16+16U\equiv48\pmod{64},
\]
so this sum has $2$-adic valuation exactly $4$, contradicting
$v_2(\sum a_i^4)=4s-3\geq5$.
\end{proof}
\begin{corollary}
There is no rational antipodal $7$-design with $34$ points for $\mu$.
\end{corollary}
\begin{proof}
Such a design would supply a rational solution of \eqref{eq:moments},
contrary to Theorem~\ref{thm:main}.
\end{proof}
\begin{remark}
The theorem is stronger than the design statement: it does not use node
distinctness, nonzero nodes, or the restrictions $|x_i|<1$.
\end{remark}
\section*{Data and computational verification}
The proof is self-contained and does not depend on computation. Accompanying
source files include a short exhaustive check of the residue alternatives and
independent exact checks used during discovery. These are supplementary
checks, not proof premises.
\bibliographystyle{amsplain}
\bibliography{references}
\end{document}
